\documentclass[anonymous, sigconf]{acmart}
\settopmatter{printacmref=false}
%%
%% \BibTeX command to typeset BibTeX logo in the docs
\AtBeginDocument{%
  \providecommand\BibTeX{{%
    Bib\TeX}}}

% %% Rights management information.  This information is sent to you
% %% when you complete the rights form.  These commands have SAMPLE
% %% values in them; it is your responsibility as an author to replace
% %% the commands and values with those provided to you when you
% %% complete the rights form.
% \setcopyright{acmlicensed}
% \copyrightyear{2018}
% \acmYear{2018}
\acmDOI{}
%% These commands are for a PROCEEDINGS abstract or paper.
\acmConference[ASP-DAC '27]
{32nd Asia and South Pacific Design Automation Conference}
{January 25--28, 2027}
{Tokyo, Japan}
% %
% %  Uncomment \acmBooktitle if the title of the proceedings is different
% %  from ``Proceedings of ...''!
% %
% %\acmBooktitle{Woodstock '18: ACM Symposium on Neural Gaze Detection,
% %  June 03--05, 2018, Woodstock, NY}
\acmISBN{}


%%
%% Submission ID.
%% Use this when submitting an article to a sponsored event. You'll
%% receive a unique submission ID from the organizers
%% of the event, and this ID should be used as the parameter to this command.
%%\acmSubmissionID{123-A56-BU3}

%%
%% For managing citations, it is recommended to use bibliography
%% files in BibTeX format.
%%
%% You can then either use BibTeX with the ACM-Reference-Format style,
%% or BibLaTeX with the acmnumeric or acmauthoryear sytles, that include
%% support for advanced citation of software artefact from the
%% biblatex-software package, also separately available on CTAN.
%%
%% Look at the sample-*-biblatex.tex files for templates showcasing
%% the biblatex styles.
%%
%%
%% The majority of ACM publications use numbered citations and
%% references.  The command \citestyle{authoryear} switches to the
%% "author year" style.
%%
%% If you are preparing content for an event
%% sponsored by ACM SIGGRAPH, you must use the "author year" style of
%% citations and references.
%% Uncommenting
%% the next command will enable that style.
%%\citestyle{acmauthoryear}


%%
%% end of the preamble, start of the body of the document source.
\begin{document}

%%
%% The "title" command has an optional parameter,
%% allowing the author to define a "short title" to be used in page headers.
\title{ZhuLong: Execution-Grounded LLM Agent for EDA Scripting with Offline API Self-Exploration}

%%
%% The "author" command and its associated commands are used to define
%% the authors and their affiliations.
%% Of note is the shared affiliation of the first two authors, and the
%% "authornote" and "authornotemark" commands
%% used to denote shared contribution to the research.
\author{Ben Trovato}
\authornote{Both authors contributed equally to this research.}
\email{trovato@corporation.com}
\orcid{1234-5678-9012}
\author{G.K.M. Tobin}
\correspondingauthor
\authornotemark[1]
\email{webmaster@marysville-ohio.com}
\affiliation{%
  \institution{Institute for Clarity in Documentation}
  \city{Dublin}
  \state{Ohio}
  \country{USA}
}

\author{Lars Th{\o}rv{\"a}ld}
\affiliation{%
  \institution{The Th{\o}rv{\"a}ld Group}
  \city{Hekla}
  \country{Iceland}}
\email{larst@affiliation.org}

\author{Valerie B\'eranger}
\affiliation{%
  \institution{Inria Paris-Rocquencourt}
  \city{Rocquencourt}
  \country{France}
}

\author{Aparna Patel}
\affiliation{%
 \institution{Rajiv Gandhi University}
 \city{Doimukh}
 \state{Arunachal Pradesh}
 \country{India}}

\author{Huifen Chan}
\affiliation{%
  \institution{Tsinghua University}
  \city{Haidian Qu}
  \state{Beijing Shi}
  \country{China}}

\author{Charles Palmer}
\affiliation{%
  \institution{Palmer Research Laboratories}
  \city{San Antonio}
  \state{Texas}
  \country{USA}}
\email{cpalmer@prl.com}

\author{John Smith}
\affiliation{%
  \institution{The Th{\o}rv{\"a}ld Group}
  \city{Hekla}
  \country{Iceland}}
\email{jsmith@affiliation.org}

\author{Julius P. Kumquat}
\correspondingauthor
\affiliation{%
  \institution{The Kumquat Consortium}
  \city{New York}
  \country{USA}}
\email{jpkumquat@consortium.net}

%%
%% By default, the full list of authors will be used in the page
%% headers. Often, this list is too long, and will overlap
%% other information printed in the page headers. This command allows
%% the author to define a more concise list
%% of authors' names for this purpose.
\renewcommand{\shortauthors}{Trovato et al.}

%%
%% The abstract is a short summary of the work to be presented in the
%% article.
\begin{abstract}
    EDA scripting with tool-specific, often undocumented APIs remains a long-tail 
    bottleneck that existing LLMs fail to address. This paper presents 
    \textbf{ZhuLong}, an execution-grounded LLM coding agent for PyAether and 
    SKILL that combines API retrieval, documentation inspection, and sandbox 
    execution via unified MCP tools, augmented by an offline \textbf{API 
    self-exploration mechanism} that infers undocumented API behaviors through 
    counterfactual experimentation.
    % 使用工具特定且常无文档的EDA API进行脚本编写，仍是现有LLM无法解决的
    % 长尾瓶颈。本文提出\textbf{ZhuLong}，一个面向PyAether和SKILL的执行驱动
    % LLM代码智能体，通过统一的MCP工具组合API检索、文档查阅和沙盒执行，
    % 并辅以离线\textbf{API自探索机制}，通过反事实实验推断未文档化的API行为。
    
    We evaluate ZhuLong on \textbf{EDA-Eval-PyAether}, a benchmark of 158
    real-world tasks with assertion-based execution, where the complete system
    achieves \textbf{78.5\%} Pass@1 in the commercial Empyrean Aether
    environment, substantially outperforming a pure LLM baseline (\textbf{23.6\%}). Ablation studies identify sandbox execution as the dominant
    performance driver (41.2 pp drop when removed), with the self-exploration
    mechanism contributing an additional 3.2 pp accuracy gain and a 22.1\%
    reduction in per-task tool calls. On 20 interactive tasks involving unsaved
    layouts and schematics, ZhuLong achieves 60.0\% Pass@1 for PyAether and
    50.0\% for SKILL. 
    % 我们在\textbf{EDA-Eval-PyAether}上评估ZhuLong，该基准包含158个真实任务，
    % 采用基于断言的执行验证，完整系统在商业Empyrean Aether环境中达到
    % \textbf{78.5\%}的Pass@1，显著优于纯LLM基线（\textbf{23.6\%}）。消融实验表明，沙盒执行是性能提升的主导因素
    % （移除后下降41.2个百分点），自探索机制在贡献额外3.2个百分点准确率提升的
    % 同时，将每任务工具调用次数减少22.1\%。在20个涉及未保存版图和原理图的
    % 交互任务中，ZhuLong在PyAether上达到60.0\% Pass@1，在SKILL上达到
    % 50.0\% Pass@1。
\end{abstract}

\keywords{LLM Agents for EDA, EDA Scripting, Code Generation, Sandbox Execution, API Self-Exploration}

% %%
% %% The code below is generated by the tool at http://dl.acm.org/ccs.cfm.
% %% Please copy and paste the code instead of the example below.
% %%
% \begin{CCSXML}
% <ccs2012>
%  <concept>
%   <concept_id>00000000.0000000.0000000</concept_id>
%   <concept_desc>Do Not Use This Code, Generate the Correct Terms for Your Paper</concept_desc>
%   <concept_significance>500</concept_significance>
%  </concept>
%  <concept>
%   <concept_id>00000000.00000000.00000000</concept_id>
%   <concept_desc>Do Not Use This Code, Generate the Correct Terms for Your Paper</concept_desc>
%   <concept_significance>300</concept_significance>
%  </concept>
%  <concept>
%   <concept_id>00000000.00000000.00000000</concept_id>
%   <concept_desc>Do Not Use This Code, Generate the Correct Terms for Your Paper</concept_desc>
%   <concept_significance>100</concept_significance>
%  </concept>
%  <concept>
%   <concept_id>00000000.00000000.00000000</concept_id>
%   <concept_desc>Do Not Use This Code, Generate the Correct Terms for Your Paper</concept_desc>
%   <concept_significance>100</concept_significance>
%  </concept>
% </ccs2012>
% \end{CCSXML}

% \ccsdesc[500]{Do Not Use This Code~Generate the Correct Terms for Your Paper}
% \ccsdesc[300]{Do Not Use This Code~Generate the Correct Terms for Your Paper}
% \ccsdesc{Do Not Use This Code~Generate the Correct Terms for Your Paper}
% \ccsdesc[100]{Do Not Use This Code~Generate the Correct Terms for Your Paper}

% \received{20 February 2007}
% \received[revised]{12 March 2009}
% \received[accepted]{5 June 2009}

%%
%% This command processes the author and affiliation and title
%% information and builds the first part of the formatted document.
\maketitle

\input{ASPDAC2027/0_abstract}
\input{ASPDAC2027/1_intro}
\input{ASPDAC2027/2_background}
\input{ASPDAC2027/3_design}
\input{ASPDAC2027/4_self_doc}
\input{ASPDAC2027/5_bench}
\input{ASPDAC2027/6_exp}
\input{ASPDAC2027/7_discussion}
\input{ASPDAC2027/8_conclusion}

%%
%% The next two lines define the bibliography style to be used, and
%% the bibliography file.
\bibliographystyle{ACM-Reference-Format}
\bibliography{reference}

\end{document}
\endinput
%%
%% End of file `sigconf.tex'.
